\clearpage
\fancyhead[L]{\small ZWEITER KOMET VON 1805.}
\hypertarget{native-275-lower}{}
\begin{center}
ZWEITER KOMET VON 1805.

\rule{0.6\linewidth}{0.4pt}

Monatliche Correspondenz. Band XIV. S. 75–86. Juli 1806.

\rule{0.6\linewidth}{0.4pt}

\textsc{Gauss an von Zach.}
\end{center}
\begin{flushright}
Braunschweig, 20. Mai 1806.
\end{flushright}

\par\indent --- --- --- Es ist schon eine ganze Weile her, dass Sie mir freundlicherweise die Originalbeobachtungen aus Marseille zu den beiden im vergangenen Jahr erschienenen Kometen handschriftlich mitgeteilt haben. Gleich nachdem ich sie erhalten hatte, nahm ich die Berechnung des zweiten Kometen wieder auf und hätte Ihnen die Ergebnisse daher schon vor zwei Monaten schicken können. Dass ich es bisher noch nicht getan habe, lag keineswegs an Nachlässigkeit. Ich hatte mir vorgenommen, meine Untersuchung auszuweiten, und gerade diese Arbeit wurde durch verschiedene andere Verpflichtungen verzögert und unterbrochen. Jetzt zögere ich jedoch nicht länger, Ihnen meine bisherigen Untersuchungen vorzulegen. Sie haben zwar nicht zu dem Ziel geführt, das ich anfangs erreichen wollte und zu erreichen hoffte, werden aber auch so schon Anlass zu mancher durchaus interessanten Überlegung geben. Nach diesen Vorbemerkungen will ich nun der Reihe nach über meine Rechnungen berichten.

\par\indent Die Beobachtungen von Herrn~\textsc{Thulis} sind an sich alle sehr gut. Die Sternpositionen, die er mit dem parallaktischen Instrument beobachtet hat, sind jedoch größtenteils sehr fehlerhaft, besonders die Rektaszensionen. So ist etwa die vom 29. November um 25$^{\mathrm{s}}$ in Zeit zu klein, die vom 6. Dezember um 30$^{\mathrm{s}}$ und die vom 8. Dezember um 38$^{\mathrm{s}}$. Auch die Deklinationen müssen zum Teil um mehrere Minuten geändert werden. Es ist daher kein Wunder, dass bei vielen Sternen vermerkt wurde, sie seien nicht aufzufinden, obwohl ich den weitaus größten Teil reduzieren konnte. Ich habe alle Beobachtungen neu reduziert, außer denen vom 11. und 22. November (für die ich
\SourceLeaf{276}%
die Sterne nirgends gefunden habe) und der vom 5. Dezember, deren Stern zwar bei \textsc{Bode} (64 Piscium), aber nirgends in der Hist.\ cél. vorkommt. Zur Identifizierung der Sterne habe ich daher noch Folgendes hinzuzufügen:

\begin{enumerate}
\item[1)] Die Sterne vom 15. und 16. November sind ein und derselbe Stern. Hist.\ cél., S.\ 477 (die Conn.\ des tems an XI besitze ich nicht).
\item[2)] Der Stern vom 23. November steht in der Hist.\ cél., S.\ 20.
\item[3)] Der Stern am 3. Dezember ist 113 Piscium bei \textsc{Bode} und steht in der Hist.\ cél., S.\ 37 und 200.
\item[4)] Am 7. Dezember ist es 374 Aquarii bei \textsc{Bode}, 981 bei \textsc{Mayer}.
\item[5)] Am 8. Dezember ist es 359 Aquarii bei \textsc{Bode}, derselbe Stern, mit dem Dr.\ \textsc{Olbers} und ich den Kometen an diesem Abend verglichen haben.
\end{enumerate}

{\rightskip=0pt plus 1em
\par\indent Nachdem ich die Beobachtungen nun sorgfältig reduziert hatte, verbesserte ich damit meine parabolischen Elemente. Dabei zeigte sich, dass sich die drei zugrunde gelegten Positionen in der Parabel nicht ganz so genau darstellen ließen, wie man es sich hätte wünschen können. Deshalb mussten die Differenzen so gut wie möglich auf alle drei Beobachtungen verteilt werden. Das habe ich mit größter Sorgfalt nach einer eigenen Methode getan, die ich auch in mein Werk aufnehmen möchte. So ergaben sich die folgenden neuen parabolischen Elemente des zweiten Kometen. Seeberger Zeit. Die Bewegung ist rechtläufig.\par}



\begin{center}\begingroup
\fontsize{9}{11}\selectfont
\setlength{\tabcolsep}{2pt}\renewcommand{\arraystretch}{1.12}
\begin{tabular}{lr}
Durchgang durch das Perihel, 1805 Dez. 31 & $6^{\mathrm{u}}\,54^{\mathrm{m}}\,34^{\mathrm{s}}$ \\
Länge des Perihels .\ .\ .\ .\ .\ . & $109^{\circ}\,21'\,50''$ \\
Länge des aufsteigenden Knotens .\ . & $250\;33\;35$ \\
Neigung der Bahn .\ .\ .\ .\ .\ .\ . & $16^{\circ}\,30'\,32''$ \\
Logarithmus des kleinsten Abstands .\ . & 9.9503300 \\
Logarithmus der mittl. tägl. Bewegung . & 0.0346333 \\
\end{tabular}\endgroup\end{center}



\par\indent Hier folgt der Vergleich dieser Elemente mit sämtlichen Beobachtungen:



\begin{center}\begingroup
\fontsize{8}{10}\selectfont
\setlength{\tabcolsep}{2pt}\renewcommand{\arraystretch}{1.12}
\begin{tabular}{|c|c|c|c|c|c|l|}
\hline
& & \shortstack{Berechnete\\Rektaszension} & \shortstack{Berechnete\\Deklination} & \multicolumn{2}{c|}{Differenz in} & \\
1805 & Mittlere Zeit & & & \shortstack{Rektasz.\\R---B} & \shortstack{Deklination\\R---B} & Beobachter \\
\hline
Nov.\ 15 & $10\,28^{\mathrm{m}}\,13^{\mathrm{s}}$ & $14^{\circ}\,33'\,0''$ & $39^{\circ}\,18'\,41''$ N. & $-2'\,30''$ & $+0'\,26''$ & Mr.\ Thulis \\
16 & $7\;\;17\;\;54$ & $14\;\;12\;\;7$ & $38\;\;58\;\;53$ & $-1\;\;58$ & $+0\;\;46$ & --- --- \\
19 & $7\;\;11\;\;17$ & $12\;\;55\;\;12$ & $37\;\;34\;\;27$ & $-0\;\;4$ & $+1\;\;7$ & --- --- \\
20 & $7\;\;3\;\;15$ & $12\;\;28\;\;23$ & $36\;\;59\;\;41$ & $-0\;\;23$ & $-0\;\;32$ & --- --- \\
21 & $9\;\;36\;\;1$ & $11\;\;57\;\;5$ & $36\;\;16\;\;16$ & $+1\;\;13$ & $-0\;\;20$ & --- --- \\
23 & $7\;\;23\;\;15$ & $11\;\;1\;\;15$ & $34\;\;46\;\;43$ & $+1\;\;0$ & $+0\;\;40$ & --- --- \\
29 & $6\;\;39\;\;44$ & $7\;\;24\;\;39$ & $26\;\;47\;\;49$ & $+1\;\;32$ & $+0\;\;1$ & --- --- \\
Dec.\ 2 & $5\;\;34\;\;9$ & $4\;\;48\;\;43$ & $18\;\;58\;\;18$ & $+1\;\;35$ & $-0\;\;52$ & Dr.\ Olbers \\
2 & $9\;\;39\;\;0$ & $4\;\;36\;\;3$ & $18\;\;20\;\;18$ & $+1\;\;47$ & $-0\;\;3$ & Mr.\ Thulis \\
3 & $5\;\;48\;\;39$ & $3\;\;39\;\;51$ & $15\;\;4\;\;52$ & $+0\;\;45$ & $-2\;\;5$ & Dr.\ Olbers \\
3 & $6\;\;30\;\;37$ & $3\;\;36\;\;52$ & $14\;\;55\;\;11$ & $+1\;\;10$ & $-0\;\;26$ & Mr.\ Thulis \\
4 & $6\;\;31\;\;5$ & $2\;\;16\;\;36$ & $10\;\;6\;\;29$ N. & $+2\;\;12$ & $-1\;\;12$ & --- --- \\
6 & $7\;\;16\;\;16$ & $358\;\;39\;\;22$ & $3\;\;39\;\;32$ S. & $+0\;\;22$ & $+0\;\;48$ & --- --- \\
7 & $7\;\;35\;\;21$ & $356\;\;9\;\;21$ & $13\;\;4\;\;56$ & $0\;\;0$ & $+0\;\;42$ & --- --- \\
8 & $5\;\;27\;\;22$ & $353\;\;20\;\;30$ & $22\;\;56\;\;24$ & $-1\;\;10$ & $+2\;\;44$ & Dr.\ Olbers \\
8 & $6\;\;46\;\;23$ & $353\;\;8\;\;33$ & $23\;\;34\;\;40$ & $-0\;\;42$ & $+2\;\;18$ & --- --- \\
8 & $6\;\;55\;\;17$ & $353\;\;7\;\;44$ & $23\;\;35\;\;37$ & $+0\;\;4$ & $-0\;\;47$ & Dr.\ Gauss \\
8 & & $353\;\;6\;\;21$ & $23\;\;40\;\;59$ & $-0\;\;33$ & $-0\;\;15$ & Dr.\ Maskelyne \\
8 & $6\;\;59\;\;39$ & $353\;\;3\;\;54$ & $23\;\;47\;\;34$ & $+0\;\;9$ & $-0\;\;33$ & Mr.\ Thulis \\
8 & $7\;\;58\;\;3$ & $352\;\;57\;\;34$ & $24\;\;6\;\;39$ & $-0\;\;23$ & .\quad . & Dr.\ Gauss \\
9 & $6\;\;4\;\;9$ & $349\;\;2\;\;44$ & $35\;\;23\;\;39$ & $-0\;\;16$ & .\quad . & Mr.\ Thulis \\
\hline
\end{tabular}\endgroup\end{center}



\par\indent Der Vergleich wurde mit voller Genauigkeit durchgeführt. Aberration, Nutation und Parallaxe wurden angemessen
\SourceLeaf{277}%
berücksichtigt, und die zur größeren Genauigkeit mitbeobachteten Sekundenbruchteile wurden erst am Ende der Rechnung weggelassen. [Die Zeit ist die Ortszeit der jeweiligen Beobachtungsorte.]

\par\indent Die Differenzen sind an sich zwar nicht viel größer, als Beobachtungsfehler bei Kometen gewöhnlich sein können. In ihrem Verlauf zeigt sich aber eine unverkennbare Regelmäßigkeit, die für die Güte der Beobachtungen von \textsc{Thulis} zu sprechen scheint. Bei meiner eigenen Beobachtung am 8. Dezember schimmerte der Kern tatsächlich so deutlich durch, dass sich der Komet mit ziemlicher Genauigkeit beobachten lassen musste. Es ist sehr zu bedauern, dass nur die Meridianbeobachtungen vom 8. Dezember aus Greenwich (aus den Zeitungen) und vom 9\textsuperscript{ten} aus Marseille bekannt sind. Aus dieser Regelmäßigkeit ließ sich voraussehen, dass man mit einer ohne die parabolische Hypothese bestimmten Bahn, ob sich nun eine Ellipse oder eine Hyperbel ergeben würde, zumindest eine erheblich bessere Übereinstimmung mit den Beobachtungen erreichen könnte. Ich habe diese Arbeit daher umso lieber übernommen, als ich gerade an einigen wichtigen Änderungen und Ergänzungen meiner Methode arbeitete und eine praktische Anwendung dafür wünschte. Das kam mir umso mehr entgegen, als gerade \emph{dieses} Beispiel eine Besonderheit hat, die bei den Asteroiden Ceres, Pallas und Juno nicht vorkommen kann. Hier ist nicht der Ort, näher darauf einzugehen.

\par\indent Das Ergebnis dieser Arbeit ist in der Tat in mehr als einer Hinsicht bemerkenswert. Ich erhielt die folgende \emph{elliptische} Bahn:

\begin{center}\begingroup
\fontsize{9}{11}\selectfont
\setlength{\tabcolsep}{2pt}\renewcommand{\arraystretch}{1.12}
\begin{tabular}{lr}
\multicolumn{2}{l}{Durchgang durch das Perihel} \\
\quad am 2. Januar 1806, Seeberger Zeit .\ . & $11^{\mathrm{u}}\;8^{\mathrm{m}}\;45^{\mathrm{s}}$ \\
Länge des Perihels .\ .\ .\ .\ . & $109^\circ\;30'\;2''3$ \\
Länge des aufsteig. Knotens .\ . & $251\;28\;22.5$ \\
Neigung der Bahn .\ .\ .\ .\ .\ .\ . & $12\;43\;10.0$ \\
Logarithmus der großen Halbachse .\ .\ . & $0.4505887$ \\
Mittlere tägliche Bewegung .\ .\ .\ .\ . & $748''383$ \\
Umlaufzeit .\ .\ .\ .\ .\ .\ .\ . & 1731 Tage 17 Std. \\
Exzentrizität .\ .\ .\ .\ .\ .\ .\ .\ .\ . & $0.6769242$ \\
Größter Abstand von der Sonne .\ .\ .\ . & $4.732625$ \\
Kleinster \quad---\ ---\ --- .\ .\ .\ . & $0.911786$ \\
Logarithmus des kleinsten Abstands .\ .\ . & $9.9598931$ \\
\end{tabular}\endgroup\end{center}

\Needspace{28\baselineskip}
\par\indent Diese Elemente stimmen folgendermaßen mit den Beobachtungen überein:

\begin{center}\begingroup
\fontsize{9}{11}\selectfont
\setlength{\tabcolsep}{2pt}\renewcommand{\arraystretch}{1.12}
\begin{tabular}{c|cc|cc|l}
& \shortstack{Berechnete\\Rektaszension} & \shortstack{Berechnete\\Deklination} & \multicolumn{2}{c|}{Differenz in} & \\
1805 & & & \shortstack{Rektasz.\\R---B} & \shortstack{Deklination\\R---B} & Beobachter \\
\hline
Nov. 15 & $14^\circ\;35'\;45''$ & $39^\circ\;17'\;56''$ N. & $+15''$ & $-19''$ & Mr. Thulis \\
16 & $14\;14\;20$ & $38\;56\;6$ & $+15$ & $-1$ & --- --- \\
19 & $12\;55\;23$ & $37\;34\;5$ & $+7$ & $+45$ & --- --- \\
20 & $12\;27\;59$ & $36\;59\;24$ & $-47$ & $-49$ & --- --- \\
21 & $11\;55\;59$ & $36\;15\;58$ & $+7$ & $-38$ & --- --- \\
23 & $10\;59\;33$ & $34\;46\;25$ & $-42$ & $+22$ & --- --- \\
29 & $7\;22\;34$ & $26\;47\;40$ & $-33$ & $-8$ & --- --- \\
Dec.\ 2 & $4\;47\;28$ & $18\;58\;28$ & $+20$ & $-42$ & Dr. Olbers \\
2 & $4\;34\;22$ & $18\;20\;44$ & $+6$ & $+13$ & Mr. Thulis \\
3 & $3\;38\;51$ & $15\;5\;11$ & $-15$ & $-1'\;46$ & Dr. Olbers \\
3 & $1\;35\;50$ & $14\;55\;37$ & $+8$ & $0$ & Mr. Thulis \\
4 & $2\;15\;55$ & $10\;7\;4$ N. & $+1'\;31$ & $-37$ & --- --- \\
6 & $358\;39\;13$ & $3\;38\;58$ S. & $+13$ & $+14$ & --- --- \\
7 & $356\;9\;28$ & $13\;4\;43$ & $+7$ & $+29$ & --- --- \\
8 & $353\;21\;10$ & $22\;56\;53$ & $-30$ & $+3\;13$ & Dr. Olbers \\
8 & $353\;9\;12$ & $23\;35\;19$ & $-3$ & $+2\;57$ & --- --- \\
8 & $353\;8\;10$ & $23\;36\;16$ & $+30$ & $-8$ & Dr. Gauss \\
8 & $353\;6\;55$ & $23\;41\;36$ & $+1$ & $+22$ & Dr. Maskelyne \\
8 & $353\;4\;18$ & $23\;48\;8$ & $+33$ & $+1$ & Mr. Thulis \\
8 & $352\;57\;46$ & $24\;7\;18$ & $-11$ & .\quad . & Dr. Gauss \\
9 & $349\;2\;56$ & $35\;25\;21$ & $-4$ & .\quad . & Mr. Thulis \\
\end{tabular}\endgroup\end{center}
\SourceLeaf{278}%
\par\indent Vergleicht man diese Differenzen mit den oben für die Parabel gefundenen, sieht man auf den ersten Blick, dass die Ellipse \emph{sehr viel} besser passt. Unter den Beobachtungen von \textsc{Thulis} beträgt die Differenz nur bei der Rektaszension vom 4. Dezember mehr als $50''$. An diesem Tag wurde ein Stern verwendet, den ich aus der Conn. des tems, An. XIII, entnommen habe, für den ich aber in der Hist. cél. keine Beobachtungen finden konnte. Vielleicht ist hier eine Berichtigung nötig. Dr. \textsc{Olbers} hatte seine Deklination vom 3. Dezember selbst als sehr zweifelhaft bezeichnet. Am 8. Dezember ließ sich die Rechnung mit den Deklinationen von Dr. \textsc{Olbers} nur dann in Übereinstimmung bringen, wenn man um 3 Minuten von der Meridianbeobachtung aus Greenwich abwich, mit der auch \textsc{Thulis} und ich bis auf Kleinigkeiten übereinstimmen. Auch die Deklination vom 9. Dezember stimmt nach der Umrechnung auf den scheinbaren Wert durch Berücksichtigung der Refraktion bei der Ellipse besser mit der Schätzung von \textsc{Thulis} überein als bei der Parabel.

\par\indent Daraus lassen sich nun die folgenden Schlüsse ziehen:

\par\indent Die Beobachtungen des Kometen lassen sich durch alle Ellipsen darstellen, die von der obigen Ellipse zur Parabel überleiten. Genauer gesagt lassen sich die Beobachtungen in einer Ellipse mit \emph{jeder} mittleren Entfernung über 2.82 noch immer besser darstellen als in irgendeiner Parabel. Je weiter sich die Ellipse von der obigen entfernt, desto schlechter stellt sie die Beobachtungen dar. Auch auf der anderen Seite wird man bis zu einer gewissen Grenze mit einer mittleren Entfernung unter 2.82 die Beobachtungen noch besser darstellen können als mit der Parabel, da die Differenzen zu beiden Seiten des Minimums nur allmählich zunehmen. Ich habe dazu zwar keine Rechnung angestellt, glaube aber doch, dass selbst die mittlere Entfernung 2.0 oder eine Umlaufzeit von weniger als drei Jahren noch ebenso gut mit allen Beobachtungen vereinbar ist wie die Parabel.

\par\indent Ich habe schon öfter gedacht und gesagt, dass bei den hundert bisher berechneten Kometen, bei denen man stets ohne Weiteres die parabolische Hypothese anwendet und schon zufrieden ist, wenn sie sich mit den Beobachtungen verträgt, wohl für den weitaus größten Teil gilt, dass sich aus den Beobachtungen allein nicht einmal sicher beweisen ließe, dass ihre Bahn einer Parabel sehr nahe kommen muss. Bei jedem Kometen, dessen Bahn berechnet wird, wäre es daher wünschenswert, eine Vorstellung davon zu haben, wie weit die Grenzen auseinanderliegen, innerhalb deren die Bahn liegen muss, damit die Rechnung nicht zu stark von der Beobachtung abweicht. Unser Komet, der in 24 Tagen eine geozentrische Bewegung von beinahe 80 Grad gezeigt hat, ist in dieser Hinsicht meiner Meinung nach ein sehr bemerkenswertes Beispiel dafür, wie sehr wir über die wahren Abmessungen der Bahn im Ungewissen bleiben. Allerdings räume ich gern ein, dasss sich über seine wahre Bahn schon etwas ziemlich Bestimmtes sagen lassen müsste, wenn der Komet noch bis in den Januar hinein gut beobachtet worden wäre (was in der südlichen Hemisphäre sehr gut möglich gewesen wäre). In dieser Hinsicht wäre mir schon sehr willkommen, wenn Sie mir die Originalbeobachtungen von \textsc{Thulis} vom 10. und 12. November und früher oder später eine genaue Bestimmung der am 11. November zum Vergleich verwendeten Sterne verschaffen könnten. Eine um fünf Tage verlängerte Beobachtungsdauer ist hier schon von einigem Wert.

\par\indent Bei aller Ungewissheit schien die so viel bessere Übereinstimmung der obigen Ellipse es doch einigermaßen wahrscheinlich zu machen, dass sich der Komet tatsächlich in einer Ellipse mit einer nicht allzu langen Umlaufzeit bewegt. Dazu kam die Ähnlichkeit mit dem Kometen von 1772. \textsc{Bessel} hatte zwar durch seine Rechnungen gefunden, dass die übrigen Elemente nicht besser übereinstimmen als bei den Parabeln, wenn man jeden dieser Kometen durch eine Ellipse mit 33 Jahren Umlaufzeit darstellt. Nach der obigen Untersuchung blieb es aber immer noch möglich, dass ein und derselbe Komet in diesen 33 Jahren vielleicht zwei, drei
\SourceLeaf{279}%
oder mehr Umläufe vollendet hatte. Der lichtschwache Komet war im vergangenen Jahr nur wegen seiner großen Nähe so auffällig geworden. Seine Lichtschwäche würde leicht erklären, warum er bei den übrigen Durchgängen durch das Perihel unbemerkt blieb, die vielleicht unter ungünstigen Umständen stattfanden. Meine Rechnungen dazu ergeben jedoch, dass auch dies nicht möglich sein wird. Ich finde nämlich, dass man die Beobachtungen beider Kometen nicht zugleich mit derselben kleinsten und größten Entfernung von der Sonne in Einklang bringen kann, ohne bei den übrigen Elementen (etwa der Länge des Perihels, der Bahnneigung und der Länge des Knotens) eine zu große Differenz zu erhalten.

\par\indent Hier ist die Tabelle mit zwei Sätzen von Elementen des Kometen von 1772. Der erste gilt für eine Parabel, bei der ich den kleinsten Abstand genauso groß angesetzt hatte wie in meinen (früheren) parabolischen Elementen des vorjährigen Kometen. Der zweite gilt für eine Ellipse mit demselben kleinsten und größten Abstand wie die obige Ellipse. Der Vergleich der jeweiligen Elemente mit denen des vorjährigen Kometen zeigt, dass man beim Kometen von 1772 unter der Voraussetzung, dass der größte und der kleinste Abstand dieselben sind wie beim Kometen von 1805, stets eine erheblich kleinere Länge des Perihels, eine erheblich größere Länge des Knotens und eine um $4^\circ$ bis $5^\circ$ größere Neigung findet als beim letzteren.

\begin{center}\begingroup
\fontsize{8}{10}\selectfont
\setlength{\tabcolsep}{2pt}\renewcommand{\arraystretch}{1.12}
\begin{tabular}{l|c|c}
& \multicolumn{2}{c}{Elemente des Kometen von 1772.} \\
 & parabolische & elliptische \\
\hline
Durchgang durch das Perihel & 9. Febr. 1772. $5^{\mathrm{u}}$ & 8. Febr. 1772. $1^{\mathrm{u}}$ \\
Länge des Perihels . . . . & $90^{\circ}\,17'$ & $97^{\circ}\,21'$ \\
Länge des aufsteigenden Knotens & $261\ \ 9$ & $263\ \ 24$ \\
Neigung der Bahn . . . . . & $20\ \ 28$ & $17\ \ 39$ \\
kleinster Abstand . . . . . . & $0.8918$ & $0.9118$ \\
größter Abstand . . . . . . & $\infty$ & $4.7326$ \\
\end{tabular}\endgroup\end{center}

\par\indent Diese Elemente beruhen nur auf der ersten und der letzten Beobachtung von \textsc{Messier} (zwei Beobachtungen reichen immer aus, wenn zwei der Elemente als gegeben angesehen werden). Für die mittlere Beobachtung ergab sich dann in der Länge eine Differenz von 6 Minuten bei der Parabel und von 3 Min. bei der Ellipse. Die Breite stimmte in beiden Fällen gut. Durch Verteilung ließen sich diese Differenzen jedoch etwa auf die Hälfte reduzieren, ohne dass sich die Elemente selbst dadurch \emph{stark} ändern könnten.

\par\indent Trotzdem will ich die Möglichkeit, dass beide Kometen identisch sind, noch nicht ausschließen. Die bisherigen Untersuchungen zeigen allerdings, dass sich ihre Bewegungen nicht zufriedenstellend durch ein und dieselbe Bahn darstellen lassen. Es bleibt aber die Möglichkeit, dass eine äußere Störung die Elemente erheblich verändert hat. Gibt man dem Kometen von 1772 eine Umlaufzeit von 33 Jahren, sodass er 1805 wieder ins Perihel kommen musste, findet sich zwar keine Annäherung an einen Hauptplaneten. Es ist aber weiterhin möglich, dass die Ellipse, aus deren Perihel er 1772 kam, und diejenige, in deren Perihel er 1805 ankam, ganz verschiedene Umlaufzeiten haben. Dazwischen könnte er einem Hauptplaneten einmal so nahe gekommen sein, dass dieser die erste Ellipse in die zweite verwandeln konnte, gleichviel ob der Komet in dieser Zeit einen oder mehrere Umläufe vollendet hat. Es bleibt also noch die freilich etwas langwierige Untersuchung, ob sich unter allen elliptischen Bahnen, die die Bewegungen der beiden Kometen darstellen können (und die, wie wir sicher wissen, innerhalb so weiter Grenzen liegen), zwei finden lassen, die ihren jeweiligen Kometen an irgendeinem Tag zwischen 1772 und 1805 an ein und dieselbe Position führen, in deren unmittelbarer Nähe sich am selben Tag ein Hauptplanet befand. In einem solchen Fall musste dieser Hauptplanet einen starken
\SourceLeaf{280}%
Einfluss auf die Elemente haben, der aber je nach der relativen Bewegung sehr unterschiedlich ausfallen musste. Daraus ließe sich eine Änderung der Elemente, wie auch immer sie beschaffen sein mag, leicht erklären. Dazu möchte ich künftig noch weitere Untersuchungen anstellen. --- --- ---